\subsection{Universal covering}

DEFINITION 1. --- A pointed set is a set X equipped with one of its elements, called the base point. The set X equipped with the point x is sometimes denoted (X, x). Let (X, x) and (Y, y) be pointed sets; a pointed map from (X, x) to (Y, y) is a map f from X to Y such that \(f(x) = y\) .

The notions of pointed topological space, pointed continuous map, pointed covering of a pointed topological space, etc., are defined analogously.

If \((X,x)\) and \((Y,y)\) are pointed topological spaces, the set of pointed continuous maps from \((X,x)\) to \((Y,y)\) is denoted \(\mathcal{C}((X,x);(Y,y))\).

Instead of saying “let f be a pointed map from \((X, x)\) on \((Y, y)\) ”, one often uses the following phrase: “let \(f: (X, x) \to (Y, y)\) be a pointed map”.

Let B be a topological space and let b be a point of B.

DEFINITION 2. --- A pointed covering (E,x) of (B,b) is said to be a universal covering if, for every pointed covering (E',x') of (B,b), there exists a unique morphism of coverings of B, \(f: \mathrm{E} \to \mathrm{E}'\) , such that \(f(x) = x'\) .

If (E, x) and (E', x') are universal coverings of (B, b), the unique B-morphism from (E, x) to (E', x') is an isomorphism of B-spaces.

Let E be a connected covering of B and let x be a point of the fibre \(E_{b}\) . Suppose that, for every pointed covering \((\mathrm{E}^{\prime}, x^{\prime})\) of \((\mathrm{B}, b)\) , there exists a morphism \(f: E \to E^{\prime}\) of coverings of B such that \(f(x) = x^{\prime}\) . Such a morphism f is then unique (I, p. 34, cor. 1 of prop. 11), hence \((\mathrm{E}, x)\) is a universal covering of \((\mathrm{B}, b)\) . *We shall see later (I, p. 126, cor. of prop. 3) that this is in particular the case if every covering of E is trivializable.*

PROPOSITION 1. --- Let B be a connected and locally connected topological space and let b be a point of B. Let (E, x) be a universal covering of (B, b). Then, E is a Galois covering of B and every covering of B is associated to E.

The space E is locally connected, because B is. Let \(E_{0}\) be the connected component of x in E, so that the space \((\mathrm{E}_{0}, x)\) is a pointed covering of \((\mathrm{B}, b)\) (I, p. 80, cor. 1 of prop. 6). There then exists a unique morphism of coverings \(f: E \to E_{0}\) such that \(f(x) = x\). If i denotes the canonical injection of \(E_{0}\) into E, the map \(i \circ f: E \to E\) is a morphism of coverings that maps x to x, as does the map \(Id_{E}\); since \((\mathrm{E}, x)\) is a universal covering of \((\mathrm{B}, b)\), we therefore have \(i \circ f = Id_{E}\). This implies that i is surjective, hence \(E_{0} = E\). Consequently, E is connected.

Let \(y\) be a point of \(\mathrm{E}_b\) and consider the pointed covering \((\mathrm{E},y)\) of \((\mathrm{B},b)\) ; by hypothesis there exists a unique morphism of coverings \(f\colon \mathrm{E}\to \mathrm{E}\) such that \(f(x) = y\) . The map \(s\colon \mathrm{E}\to \mathrm{E}\times_{\mathrm{B}}\mathrm{E}\) defined by \(t\mapsto (t,f(t))\) is a continuous section of the map \(\mathrm{pr}_1\colon \mathrm{E}\times_{\mathrm{B}}\mathrm{E}\to \mathrm{E}\) . By cor. 4 of I, p. 81, this shows that the covering \(\mathrm{E}\times_{\mathrm{B}}\mathrm{E}\) of E given by the map \(\mathrm{pr}_1\) is trivializable. The covering E is therefore Galois (th. 2 of I, p. 102). It then follows from I, p. 112, corollary of prop. 10, that every covering of B is associated to E.

COROLLARY. --- Let B be a locally connected topological space. Let b be a point of B and let (E, x) be a universal covering of (B, b). For a subspace A of B, the following two properties are equivalent:

\begin{enumerate}
\def\labelenumi{(\roman{enumi})}
\item
  The covering E is trivializable over A;
\item
  Every covering of B is trivializable over A.
\end{enumerate}

Property (ii) obviously implies property (i). The converse follows from the fact that every covering of B is associated to the covering E (I, p. 105, prop. 7).

